% AI-assisted (Claude, writer), 26 Nov 2026: restructure after R1-ALIGN-OUTCOME (task writer/20261126-01).
% Text marked MOVED is carried over unchanged from the 23 Nov main text unless a comment says otherwise.
% Not an author decision. No label, test, evaluation or audit output was opened to write it.
% MOVED from related.tex: the full related-work text of the 23 Nov build. The main text carries a shorter version.
\section{Extended Related Work}
\label{app:related}

\paragraph{Position.}
Noise-driven acceptance in automated research loops is documented, and statistical gates exist that control it. This paper differs from that work in four respects. (a) Our acceptance and screening rule maximises audited gain per run under a declared run price; the gate family controls a false-accept probability. (b) Screening happens before a run is spent, from a per-proposal forecast of the increment distribution. (c) The theory covers a proposer whose law depends on the ledger: a winner's curse for any adaptive proposer, and conditions under which the one-step rule is globally optimal, with a failure example. (d) We compare a learned forecast with a theory-derived index policy at matched compute. We do not claim the first measurement of the noise problem. Our measurements add audited ground truth, several nodes and two proposers under governance. Two practices from environment synthesis for agent training carry over. An environment is accepted only if a reference solution passes and a no-operation agent fails \citep{xu2026skill2env}; we report the analogous positive control for each node, with fresh seeds, since a single trial can pass by chance; the negative control has no counterpart here, because the harness does not run an edit that changes nothing. The same work evaluates behaviour with the judge that supplied its training reward; our audits instead use seeds disjoint from those that informed any decision.

\paragraph{Measurements of noise in research loops.}
\citet{jalaian2026noise} replay accept rules on real runs, with 16 conditions each re-run 30 times and 1{,}280 logged hyperparameter trials. Deciding on one run keeps 49.9\% of truly neutral candidates. Peeking after every run gives 11.3\% false accepts at a nominal 5\% level, a Robbins normal-mixture e-process holds 5.3\%, and a 20-run calibration phase holds the nominal level. The study has no learned predictor, no cost-aware Bayes rule and no screening before execution. \citet{grover2026lucky} uses a controlled nanoGPT setting with 60 seeds. Best-of-32 selection alone gives an apparent gain of 0.00544 in validation bits per byte, 2.7 times a preregistered threshold of 0.002. In a frozen bank of 24 candidates evaluated over 749 runs, nine were exact nulls; of the other 15, eight replicated, three failed and four were inconclusive. That study measures and confirms; it proposes no decision rule. Our C1 is supporting evidence of the same phenomenon in a different regime: bounded edits proposed by a language model, with error rates computed against fresh-seed audits of every eligible proposal.

\paragraph{Statistical acceptance gates.}
SAGE \citep{wang2026sage} gates self-evolving agents that edit a skill document. It compares candidate and incumbent per item and applies a one-sided exact binomial test to the discordant pairs, with a regression cap, on five benchmarks and four backbones. PACE \citep{shawn2026pace} is an anytime-valid paired commit gate: a betting e-process on McNemar discordant pairs that bounds each candidate's false-commit probability under optional stopping. The Statistical G\"odel Machine \citep{sgm2025} applies statistical tests with a global error budget to recursive self-modification. \citet{lonkar2026evidence} freezes a candidate, draws fresh audit samples and requires Hoeffding lower confidence bounds, non-inferiority on protected capabilities and a canary deployment; the evaluation is a synthetic Monte Carlo study, and the method has no model of the value of acceptance or of the cost of a run. The Ladder \citep{blum2015ladder} is the closest classical keep rule: it updates a record only when a submission beats it by a margin, with guarantees under adaptive submissions.

These gates share three properties. They control an error probability. They decide after the candidate has been evaluated, and several spend extra evaluations to do so. They involve no learning. Our rule answers a different question: which threshold, and which decision to run at all, gives the most audited gain per run at a declared price. A sceptical forecast raises the keep threshold, as a gate would, but a confident forecast lowers it, and a low forecast value skips the run. The two views can be combined; we do not study the combination.

\paragraph{Prediction before execution.}
FOREAGENT \citep{zheng2026foreagent} prompts a language model for pairwise preferences between complete solutions and executes the top-ranked candidate. It does not fine-tune a predictor, does not predict an increment distribution, and has no noise model or decision threshold. Its reported pairwise accuracy is modest, and lower for small open models, which is a warning for our fine-tuned arm. Beneath the Diff \citep{he2026beneath} studies the same keep-if-better edit loop, documents repetitive edits and filters similar proposals; it gives no statistical model of the keep decision. \citet{chen2026visible} show that proposers re-propose rejected candidates in the states that a loop reaches, with duplicates taking between 28\% and 63\% of valid iterations in a deterministic setting. In-context surrogates inside Bayesian optimisation predict learning curves or outcomes \citep{arxiv240416795,liu2024llambo}; they address a different decision, namely which configuration to evaluate or continue. Our screening rule is value-based. It needs the no-op probability, the mean and the spread of the increment, and it compares the resulting value with a run price. Rule-detectable duplicates are handled by a rule layer and reported as a separate arm, so that a learned gain is not credited for duplicate detection.

\paragraph{Decision-focused training and learned controllers for agents.}
The closest relative of our decision-focused arm is \citet{zhang2026calibrationcontrol}. They argue that runtime oversight should target intervention advantage, the expected utility gain from intervening instead of continuing, in place of a calibrated failure score, and they evaluate action-conditioned control on agent benchmarks that include ALFWorld. Making decision utility, not calibration, the target of the learned object therefore has precedent, and we do not claim otherwise. A proactive desktop agent learns, per user, when to intervene from engagement and dismissals \citep{shao2026deeppersonalization}; its decision is a direct yes or no from a language model with no explicit cost of interrupting, and its feedback depends on its own decisions, since steps where it stayed silent receive only a retrospective guess. Our rule carries an explicit run price, and our labels come from audits that do not depend on the gate's decision. Predict-then-optimise methods in operations research train predictors for downstream decisions in general \citep{arxiv171008005,wilder2019melding}. Our difference is the object and the rule: the predictor forecasts an audited increment distribution for a proposed experiment, and it feeds a closed-form keep and skip rule under a run price.

Several works learn a controller over a frozen language model. ARC \citep{arxiv260211574} trains a policy that selects an agent configuration (workflow, tools, budget) per query. It configures the agent, not an acceptance rule under evaluation noise. Harness-MDP \citep{yi2026harnessmdp} trains a small controller offline to choose among structural harness actions; none of its domains is machine-learning experimentation. Other neighbours fix or prompt what we derive: a fixed rate at which a language model proposes inside an evolution strategy \citep{ferreira2026centaur}, a small advisor trained on counterfactual branches \citep{jiang2026cota}, budgeted allocation of research effort \citep{yu2026primescientist}, and a prompted remedy for a judge that degrades late in a campaign \citep{ji2026rehearse}. We do not claim that a small learned controller over a language model is new.

\paragraph{Dynamic algorithm configuration.}
DAC adapts the parameters of an algorithm while it runs, with a policy learned across instances \citep{biedenkapp2020dac,adriaensen2022dac,eimer2021dacbench}. We use its vocabulary with a specific mapping (Section~\ref{sec:problem}): the configured parameters are two thresholds and a verification count, and instances are tuples of node, seed, manager and budget. Two points differ from the usual DAC setting. The learned object is a forecaster, and the policy is derived from it. And we state when the dynamic problem reduces to one-step decisions, as a conditional theorem with synthetic controls; we do not present the reduction as an empirical finding.

\paragraph{Selection under noise.}
The formulas behind our rule are classical. The Bayes threshold against a standard and the economics of sampling appear in ranking and selection \citep{xie2013sequential,chick2012sequential}, the value of one more measurement in the knowledge gradient \citep{frazier2008knowledge}, repeated evaluation in racing \citep{birattari2002frace,lopezibanez2016irace}, and sample allocation in best-arm identification \citep{audibert2010bestarm,kaufmann2016bestarm}. The optimizer's curse \citep{smith2006optimizer} explains why a selected estimate is biased. We present these formulas as instantiations. What those models assume is a fixed set of alternatives. The loop departs from them in three ways: the proposer's law depends on the ledger, a discarded candidate restores the incumbent, and the selection bias of the record feeds back into what the proposer sees.
